\section{Functional and Performance Requirements}
\label{sec:req-functional}

The FR labels identify control functions, and the PR labels identify performance tests in the results chapter.

MCU pin assignments are listed in \cref{sec:intro-controller-basics}. On the original prototype, pin PA10 controls switching bursts, PA1 senses secondary bulk voltage, PA0 senses primary voltage, and PA2 senses half-bridge current. DMA transfers ADC readings to memory, and UART carries commands and status.

\begin{table}[H]
\centering
\caption{Functional requirements.}
\label{tab:req-functional}
\begin{tabularx}{\textwidth}{@{}L{0.12\textwidth}Y@{}}
\toprule
ID & Required behavior\\
\midrule
FR-01 & Start through a controlled sequence that permits the isolated secondary feedback path to become valid before normal regulation.\\
FR-02 & Generate the fixed carrier and regulate average transferred power through PA10 burst enable.\\
FR-03 & Use the PA1 secondary-voltage measurement for final closed-loop operation.\\
FR-04 & Reduce transferred energy at no load and respond to load removal without uncontrolled secondary-voltage accumulation.\\
FR-05 & Recover from load attachment and load changes without an unintended persistent fault.\\
FR-06 & Apply the configured secondary-overvoltage response: eligible no-load soft holdoff, or a latched PA10 inhibit when the hard-fault criteria are met.\\
FR-07 & Acquire PA0, PA1, and PA2 through DMA during active control.\\
FR-08 & Provide UART access to version, mode, settings, calibration, status, and fault clearing.\\
\bottomrule
\end{tabularx}
\end{table}

\begin{table}[H]
\centering
\caption{Performance tests and measurement conditions.}
\label{tab:req-performance}
\begin{tabularx}{\textwidth}{@{}L{0.12\textwidth}Y@{}}
\toprule
ID & Measurement or test\\
\midrule
PR-01 & Assess the \SI{65}{\watt} continuous minimum and earlier \SI{75}{\watt} validation milestone against the \SI{100}{\watt} continuous target.\\
PR-02 & Identify the AC input conditions used in the tests.\\
PR-03 & Identify the output settings actually obtained under load.\\
PR-04 & Report mean VBUS, ripple observations, and available load-transient excursions.\\
PR-05 & Assess observed no-load voltage against the configured control and protection thresholds where the records permit.\\
PR-06 & Report PF and input-current distortion at selected operating points.\\
PR-07 & Calculate efficiency from input and output power measured under the same conditions, with the included circuit stages stated.\\
PR-08 & Report temperature together with load, run time, ambient condition, and airflow information.\\
PR-09 & Record the assembly, firmware, operating conditions, and source measurements for each numerical result.\\
PR-10 & Identify the limiting mechanism in any reported overload or current-limit test.\\
\bottomrule
\end{tabularx}
\end{table}

Chapter~\ref{ch:results} compares implementation with measured electrical behavior. Section~\ref{sec:revised-plan-criteria} lists the additional V2/V3 design and test tasks under AP labels.
